\documentclass[exercice,a4]{thedocument}%
\usepackage[french]{babel}%
\usepackage{pstricks}%
\usepackage{array}%
\usepackage[np]{numprint}%
\usepackage{eurosym}%
\usepackage{amsmath}%
\startdatakeys
    \source{22GENMATMEAG2}
    \BO{2026}
    \cycle{4}
    \competencesDuSocle{CA,RA}
    \objectifsApprentissage{C4B3ca,C4B1ed,C4B1bf}
    \motsCles{Thalès,droites parallèles,volume,cylindre,aire,agrandissement}
\enddatakeys
\begin{document}%

\medskip

La tour de la Vade est un monument de Carcassonne.

\begin{minipage}{0.6\linewidth}
Représentation schématique de la tour de la Vade (sans le toit).
\end{minipage}\hfill
\begin{minipage}{0.38\linewidth}
\begin{center}
\psset{unit=1cm}
\begin{pspicture}(-1,-2.3)(1,0.4)
\psellipse(0,0)(1,0.4)
\psline(-1,0)(-1,-1.8)\psline(1,0)(1,-1.8)
\psellipticarc(0,-1.8)(1,0.4){180}{360}
\psellipticarc[linestyle=dashed](0,-1.8)(1,0.4){0}{180}
\end{pspicture}
\end{center}
\end{minipage}

\begin{enumerate}
\item Afin de déterminer la hauteur de cette tour, Romane et Maël se sont positionnés comme indiqué sur la figure ci-dessous, et ont effectué plusieurs mesures.

L'oeil de Maël est au point M ; le segment [FG] représente Romane.

La figure n'est pas à l'échelle.

\begin{center}
\psset{unit=1cm}
\begin{pspicture}(9,3.6)
\pspolygon(1,0.4)(8.4,0.4)(8.4,3.1)(7.1,3.1)%MCDA
\psframe[fillstyle=solid,fillcolor=lightgray](7.1,0.4)(8.4,3.1)
\psline[linewidth=1.5pt](3.3,0.4)(3.3,1.45)%GF
\uput[u](7.1,3.1){A} \uput[d](7.1,0.4){B} \uput[dr](8.4,0.4){C} \uput[u](8.4,3.1){D}
\uput[ul](3.3,1.45){F} \uput[d](3.3,0.4){G} \uput[ul](1,0.4){M}
\uput[d](1,0.4){Maël}\rput{90}(3.5,0.9){\small Romane}\rput(7.75,1.75){Tour}
\psframe(3.3,0.4)(3.1,0.6)\psframe(7.1,0.4)(6.9,0.6)
\end{pspicture}
\end{center}

Les points M, F et A ainsi que les points M, G et B sont alignés.

Romane et Maël ont mesuré : MG $= 3$~m

\phantom{Romane et Maël ont mesuré : }FG $= 1,4$ m

\phantom{Romane et Maël ont mesuré : }GB $= 51$ m

	\begin{enumerate}
		\item Montrer que les droites (FG) et (AB) sont parallèles.
		\item Vérifier que la hauteur AB de la tour est de $25,2$~m.
	\end{enumerate}
\item La tour de la Vade a une base circulaire de diamètre proche de $14$~m.

Montrer que son volume est d'environ \np{3880}~m$^3$.
\item Romane a acheté une maquette de cette tour à l'échelle $\dfrac{1}{20}$.

Quel est le volume de cette maquette, réduction de la tour de la Vade ?
\item La tour doit être entretenue ; il faut passer un traitement contre la moisissure sur toute sa
surface. Le coût du traitement est de $39$~\euro{} par m$^2$.

Combien va coûter le traitement de la tour ?

On rappelle que :

\begin{tabular}{|l|}\hline
Aire latérale d'un cylindre $= 2 \times \pi \times \text{rayon} \times \text{hauteur}$\\
Volume d'un cylindre $= \text{aire de la base} \times  \text{hauteur}$\\ \hline
\end{tabular}
\end{enumerate}

\showthesource
\end{document}
